% AA-71
\chapter{Wzór wymiarowy Weyla}
\label{aa:ch:weyl-dimension}

Wzór charakterów Weyla daje wszystkie krotności wag, ale jego licznik i mianownik znikają po podstawieniu $e^\mu=1$. W tym rozdziale obliczymy dokładny rząd znikania obu wyrażeń. Wykorzystamy symetrię alternantów, która pozwala sprowadzić również licznik do tożsamości mianownika.

Zachowujemy założenia AA-70: $K$ jest algebraicznie domkniętym ciałem charakterystyki zero, $\mathfrak g$ niezerową skończenie wymiarową półprostą algebrą Liego nad $K$, $\mathfrak h$ jej maksymalną podalgebrą toralną, $\Phi^+$ zbiorem korzeni dodatnich względem bazy $(\alpha_1,\ldots,\alpha_r)$, $P$ kratą wag, $P_+$ zbiorem wag dominujących, $W$ grupą Weyla, $\rho=\sum_i\omega_i=\frac12\sum_{\alpha\in\Phi^+}\alpha$. Przez $(\cdot,\cdot)$ oznaczamy dodatnio określony iloczyn skalarny w rzeczywistym modelu $E=P\otimes_{\mathbb Z}\mathbb R$ pochodzący z formy Killinga (AA-60 i AA-70). Krotności wag są liczbami całkowitymi, a obliczenia granic wykonujemy w $\mathbb R$; nie zakładamy istnienia włożenia $K$ do $\mathbb R$ ani do $\mathbb C$.

\section{Specjalizacja alternantów}

\begin{definicja}
\label{aa:def:weyl-exponential-specialization}
Dla $v\in E$ i $t\in\mathbb R$ określamy specjalizację
$\mathcal E_{v,t}:\mathbb R[P]\to\mathbb R$ przez liniowe rozszerzenie wzoru
$\mathcal E_{v,t}(e^\mu)=\exp(t(\mu,v))$.
Przez $N$ oznaczamy liczbę korzeni dodatnich, czyli $N=|\Phi^+|$.
\end{definicja}
Odwzorowanie $\mathcal E_{v,t}$ jest homomorfizmem pierścieni, gdyż addytywność iloczynu skalarnego daje
$\exp(t(\mu+\nu,v))=\exp(t(\mu,v))\exp(t(\nu,v))$.
Dla $t=0$ specjalizacja jest homomorfizmem sumującym współczynniki. W szczególności
$\mathcal E_{v,0}(\operatorname{ch}L(\lambda))=\dim_K L(\lambda)$.
Dla ustalonego elementu pierścienia funkcja $t\mapsto\mathcal E_{v,t}(F)$ jest skończoną sumą funkcji wykładniczych, a więc jest analityczna.

\begin{lemat}[Symetria alternantów]
\label{aa:lem:weyl-alternant-specialization}
Przy powyższych założeniach dla $\mu,\nu\in P$ i $t\in\mathbb R$ zachodzi
\[
 \mathcal E_{\nu,t}(A_\mu)=\mathcal E_{\mu,t}(A_\nu).
\]
W szczególności
\[
 \mathcal E_{\rho,t}(A_\mu)=
 \exp(t(\rho,\mu))
 \prod_{\alpha\in\Phi^+}(1-\exp(-t(\alpha,\mu))).
\]
\end{lemat}
\begin{proof}
Grupa Weyla działa przez izometrie, zatem $(w\mu,\nu)=(\mu,w^{-1}\nu)=(w^{-1}\nu,\mu)$. Ponadto $\epsilon(w^{-1})=\epsilon(w)$, gdyż $\epsilon(w)=\det w$ należy do $\{-1,1\}$. Podstawienie $u=w^{-1}$ jest bijekcją skończonej grupy $W$, więc
\[
 \begin{aligned}
 \mathcal E_{\nu,t}(A_\mu)
 &=\sum_{w\in W}\epsilon(w)\exp(t(w\mu,\nu))\\
 &=\sum_{u\in W}\epsilon(u)\exp(t(u\nu,\mu))
 =\mathcal E_{\mu,t}(A_\nu).
 \end{aligned}
\]
Bierzemy $\nu=\rho$ i stosujemy Twierdzenie~\ref{aa:tw:weyl-denominator} do $A_\rho$. Ponieważ specjalizacja zachowuje iloczyny, otrzymujemy drugą równość.
\end{proof}

\begin{lemat}[Rząd znikania]
\label{aa:lem:weyl-alternant-order}
Przy powyższych założeniach niech $\mu\in P$ będzie ściśle dominująca. Funkcja $t\mapsto\mathcal E_{\rho,t}(A_\mu)$ ma w zerze zero rzędu dokładnie $N$. Zachodzi
\[
 \lim_{t\to0}\frac{\mathcal E_{\rho,t}(A_\mu)}{t^N}
 =\prod_{\alpha\in\Phi^+}(\mu,\alpha)>0.
\]
Równoważnie, dla $0\le k<N$ mamy
$\sum_{w\in W}\epsilon(w)(w\mu,\rho)^k=0$, a
\[
 \sum_{w\in W}\epsilon(w)(w\mu,\rho)^N
 =N!\prod_{\alpha\in\Phi^+}(\mu,\alpha).
\]
\end{lemat}
\begin{proof}
Dla korzenia prostego $\alpha_i$ mamy $(\mu,\alpha_i)=\frac12(\alpha_i,\alpha_i)\mu(h_i)>0$. Każdy korzeń dodatni jest niezerową kombinacją korzeni prostych o nieujemnych współczynnikach, więc $b_\alpha=(\mu,\alpha)>0$ dla każdego $\alpha\in\Phi^+$. Z rozwinięcia funkcji wykładniczej w zerze
$1-\exp(-tb_\alpha)=tb_\alpha+O(t^2)$.
Zatem każdy czynnik ma zero proste w zerze. Po podzieleniu przez $t$ każdy z $N$ czynników dąży do $b_\alpha$, a czynnik $\exp(t(\rho,\mu))$ dąży do jeden. Iloczyn skończenie wielu granic daje żądaną granicę, która jest niezerowa.

Dokładniej, analityczność każdego czynnika oznacza, że można zapisać $1-\exp(-tb_\alpha)=t f_\alpha(t)$, gdzie $f_\alpha$ jest analityczna w otoczeniu zera i $f_\alpha(0)=b_\alpha$. Z Lematu~\ref{aa:lem:weyl-alternant-specialization} wynika więc
$\mathcal E_{\rho,t}(A_\mu)=t^N\exp(t(\rho,\mu))\prod_\alpha f_\alpha(t)$.
Czynnik mnożący $t^N$ ma w zerze wartość $\prod_\alpha b_\alpha\ne0$. Stąd współczynniki Taylora rzędów $0,\ldots,N-1$ znikają, a współczynnik rzędu $N$ jest równy temu iloczynowi.

Z drugiej strony rozwinięcie skończonej sumy definiującej alternant daje współczynnik przy $t^k$ równy
$\frac1{k!}\sum_{w\in W}\epsilon(w)(w\mu,\rho)^k$.
Porównanie obu rozwinięć daje ostatnie dwie tożsamości.
\end{proof}

\section{Wzór i jego konsekwencje}

\begin{twierdzenie}[Wzór wymiarowy Weyla]
\label{aa:tw:weyl-dimension}
Przy powyższych założeniach dla każdej $\lambda\in P_+$ zachodzi
\[
 \dim_K L(\lambda)=
 \prod_{\alpha\in\Phi^+}
 \frac{(\lambda+\rho,\alpha)}{(\rho,\alpha)}
 =\prod_{\alpha\in\Phi^+}
 \frac{(\lambda+\rho)(h_\alpha)}{\rho(h_\alpha)}.
\]
\end{twierdzenie}
\begin{proof}
Ze wzoru charakterów Weyla (Twierdzenie~\ref{aa:tw:weyl-character-formula}) mamy
$A_\rho\operatorname{ch}L(\lambda)=A_{\lambda+\rho}$.
Z Lematu~\ref{aa:lem:weyl-rho} obie wagi $\rho$ i $\lambda+\rho$ są ściśle dominujące. Dla $t\ne0$ każdy czynnik $1-\exp(-t(\alpha,\rho))$ jest niezerowy, ponieważ $(\alpha,\rho)>0$ i rzeczywista funkcja wykładnicza przyjmuje wartość jeden tylko dla argumentu zero. Z Lematu~\ref{aa:lem:weyl-alternant-specialization} wynika $\mathcal E_{\rho,t}(A_\rho)\ne0$. Stosujemy homomorfizm $\mathcal E_{\rho,t}$ do wzoru charakterów i dzielimy przez tę liczbę:
\[
 \mathcal E_{\rho,t}(\operatorname{ch}L(\lambda))
 =\frac{\mathcal E_{\rho,t}(A_{\lambda+\rho})/t^N}
        {\mathcal E_{\rho,t}(A_\rho)/t^N}.
\]
Lewa strona dąży dla $t\to0$ do $\dim_K L(\lambda)$, gdyż charakter jest skończoną sumą. Z Lematu~\ref{aa:lem:weyl-alternant-order} licznik i mianownik po prawej stronie dążą odpowiednio do iloczynów $\prod_\alpha(\lambda+\rho,\alpha)$ i $\prod_\alpha(\rho,\alpha)$, przy czym drugi iloczyn jest dodatni. Granica ilorazu istnieje i daje pierwszy wzór.

Dla dowolnego korzenia $\alpha$ relacja z Twierdzenia~\ref{aa:tw:lie-root-realization} daje
$\mu(h_\alpha)=2(\mu,\alpha)/(\alpha,\alpha)$ (najpierw dla korzeni, a następnie przez liniowość dla całego $E$, ponieważ korzenie rozpinają $E$). W każdym czynniku ilorazu mnożnik $2/(\alpha,\alpha)$ skraca się. Otrzymujemy drugi wzór.
\end{proof}

\begin{wniosek}
\label{aa:cor:weyl-dimension-polynomial}
Dla ustalonej $\lambda\in P_+$ i $m\in\mathbb Z_{\ge0}$ zachodzi
\[
 \dim_K L(m\lambda)=
 \prod_{\alpha\in\Phi^+}
 \left(1+m\frac{(\lambda,\alpha)}{(\rho,\alpha)}\right).
\]
Jest to wielomian w $m$ o współczynnikach wymiernych i stopniu równym liczbie korzeni dodatnich $\alpha$ takich, że $(\lambda,\alpha)>0$. W szczególności dla ściśle dominującej $\lambda$ stopień wynosi $N$, a dla $\lambda=0$ wymiar wynosi jeden.
\end{wniosek}
\begin{proof}
Podstawiamy $m\lambda$ do Twierdzenia~\ref{aa:tw:weyl-dimension} i rozdzielamy licznik każdego czynnika. Dla wagi dominującej $(\lambda,\alpha_i)\ge0$, więc również $(\lambda,\alpha)\ge0$ dla dodatnich korzeni. Czynnik jest stały i równy jeden dokładnie wtedy, gdy $(\lambda,\alpha)=0$; w pozostałych przypadkach jest wielomianem liniowym o dodatnim współczynniku wiodącym. Iloczyn ma zatem wskazany stopień i niezerowy współczynnik wiodący. Wymierność wynika z formy na kracie wag opisanej w AA-70. Dla $\lambda=0$ wszystkie czynniki są równe jeden.
\end{proof}

\section{Przykłady obliczeniowe}

\begin{przyklad}[Typy $A_1$ i $A_2$]
\label{aa:prz:weyl-dimension-a}
Dla typu $A_1$ waga $\lambda=a\omega_1$, $a\ge0$, daje $\dim L(a\omega_1)=a+1$. W typie $A_2$ dodatnie korzenie to $\alpha_1,\alpha_2,\alpha_1+\alpha_2$ i mają tę samą długość. Ich korzenie dualne mają współrzędne $(1,0),(0,1),(1,1)$ w bazie $h_1,h_2$, gdyż $h_\alpha=2t_\alpha/(\alpha,\alpha)$ i $\alpha\mapsto t_\alpha$ jest liniowe. Dla $\lambda=a\omega_1+b\omega_2$ otrzymujemy
\[
 \dim L(a\omega_1+b\omega_2)=
 \frac{(a+1)(b+1)(a+b+2)}{2}.
\]
Na przykład dla $(a,b)=(1,0),(0,1),(1,1),(2,2)$ wymiary wynoszą odpowiednio $3,3,8,27$. Pierwsze trzy wartości odpowiadają reprezentacji naturalnej, dualnej i dołączonej algebry $\mathfrak{sl}_3(K)$.
\end{przyklad}

\begin{przyklad}[Typ $B_2$]
\label{aa:prz:weyl-dimension-b2}
Niech $\alpha_1$ będzie krótkim, a $\alpha_2$ długim korzeniem prostym. Przy wspólnym przeskalowaniu iloczynu skalarnego możemy przyjąć $(\alpha_1,\alpha_1)=1$, $(\alpha_2,\alpha_2)=2$, $(\alpha_1,\alpha_2)=-1$. Współrzędne dodatnich korzeni oraz odpowiadających im elementów $h_\beta$ są następujące:
\[
 \begin{array}{c|c|c}
 \beta & (\beta,\beta)&h_\beta\\ \hline
 \alpha_1&1&h_1\\
 \alpha_2&2&h_2\\
 \alpha_1+\alpha_2&1&h_1+2h_2\\
 2\alpha_1+\alpha_2&2&h_1+h_2
 \end{array}
\]
W ostatniej kolumnie użyto wzoru
$h_\beta=\sum_i n_i\frac{(\alpha_i,\alpha_i)}{(\beta,\beta)}h_i$
dla $\beta=\sum_i n_i\alpha_i$, który wynika z liniowości $t_\beta$; skala formy skraca się w ilorazach. Ponieważ $\rho(h_i)=1$, dla $\lambda=a\omega_1+b\omega_2$ mamy
\[
 \dim L(\lambda)=
 \frac{(a+1)(b+1)(a+2b+3)(a+b+2)}{6}.
\]
Wymiary modułów o wagach $\omega_1,\omega_2,2\omega_1$ wynoszą odpowiednio $4,5,10$. Ostatnia waga to najwyższy korzeń $2\alpha_1+\alpha_2=2\omega_1$; wymiar dziesięć zgadza się z ośmioma przestrzeniami korzeniowymi i dwuwymiarową podalgebrą Cartana.
\end{przyklad}

\begin{przyklad}[Typ wyjątkowy $G_2$]
\label{aa:prz:weyl-dimension-g2}
Niech $\alpha_1$ będzie krótki, a $\alpha_2$ długi. Przy wspólnym przeskalowaniu przyjmujemy kwadraty długości $2,6$ oraz $(\alpha_1,\alpha_2)=-3$. Z listy dodatnich korzeni typu $G_2$ (AA-56) i wzoru na $h_\beta$ otrzymujemy
\[
 \begin{array}{c|c|c}
 \beta&(\beta,\beta)&h_\beta\\ \hline
 \alpha_1&2&h_1\\
 \alpha_2&6&h_2\\
 \alpha_1+\alpha_2&2&h_1+3h_2\\
 2\alpha_1+\alpha_2&2&2h_1+3h_2\\
 3\alpha_1+\alpha_2&6&h_1+h_2\\
 3\alpha_1+2\alpha_2&6&h_1+2h_2
 \end{array}
\]
Zatem
\[
 \begin{aligned}
 \dim L(a\omega_1+b\omega_2)
 &=\frac{(a+1)(b+1)(a+b+2)}{120}\\
 &\quad\cdot(a+2b+3)(a+3b+4)(2a+3b+5).
 \end{aligned}
\]
Dla $(a,b)=(1,0)$ otrzymujemy $7$, dla $(0,1)$ otrzymujemy $14$, a dla $(2,0)$ otrzymujemy $27$. Najwyższy korzeń $3\alpha_1+2\alpha_2$ ma współrzędne $(0,1)$ względem wag fundamentalnych, więc wartość $14$ jest zgodna z wymiarem reprezentacji dołączonej: dwanaście przestrzeni korzeniowych i podalgebra Cartana rzędu dwa.
\end{przyklad}

Wzór wymiarowy jest konsekwencją wzoru charakterów, ale nie odtwarza wszystkich krotności wag. Na przykład w typie $A_2$ wagi $\omega_1$ i $\omega_2$ wyznaczają nieizomorficzne moduły o tym samym wymiarze trzy.
